% page 81 of the facsimile (image p-080 of the scan)
% block 165.1 x 246.3 mm ; left margin 36.9 mm
\pgc{15.240mm}{0.000mm}{
\l{}
\l{}
\l{}
\l{\cel{51}73}
\l{}
\l{\sou{brasar}\cel{2}(\sou{trans.}) Movar, agitar por operaco. (Ex.: brasar oro}
\l{\cel{3}ed arjento en kruzelo). - FIS.}
\l{}
\l{\sou{brava}. Pronta afrontar la danjero (precipue en la kombati).-}
\l{\cel{3}DEFIRS.}
\l{}
\l{\sou{\textquotedbl{}bravur-ario\textquotedbl{}}. Ario briloza, destinita ad eminentigar la}
\l{\cel{3}kantisto. - DEF.}
\l{}
\l{\sou{brazar}. (\sou{trans.}) Unionar (du metali interdiferanta) per}
\l{\cel{3}aloyuro quan onu fuzas en fairego. - EF.}
\l{}
\l{\sou{brechio}. Roko di qua la strukturo esas fragmentoza, ek peceti}
\l{\cel{3}anguloza, diversa-kolora, aglomerita en cemento naturala.}
\l{}
\l{\sou{brecho}. Rupturo qua posibligas enirar remparo-cirkuito. - DEFIRS.}
\l{}
\l{\sou{breko}. Veturo quar-rota, generale sen-tekta, qua havas, ye}
\l{\cel{3}sua parto avana, sidili situita alte, e du benki longesale.}
\l{\cel{3}- DEFIS.}
\l{}
\l{\sou{brelano}. Hazardo-ludo, quan onu pleis per distributar tri}
\l{\cel{3}karti ad omna ludonto. - DFI.}
\l{}
\l{\sou{bremo}. (\sou{zool.}) Fisho river-aquala, qua similesas karpo, e di}
\l{\cel{3}qua la karno nutro-prizesas. - L.\cel{1}\sou{abramis brama}. - EFS.}
\l{}
\l{\sou{bretelo}.\cel{2}Bendo de ledro, quan onu pasigas sur la shultro,}
\l{\cel{3}por suportar sako, dorso-korbo, kamilo ; benda de texuro}
\l{\cel{3}elastika, bifurkoza, por subtenar la pantalono virala.-FI.}
\l{}
\l{\sou{brevo}. Reskripto pontifikala, ordinare en linguo Latina;}
\l{\cel{3}papo-sigluroza. - DEFIRS.}
\l{}
\l{\sou{breviario.}\cel{2}(\sou{religio katolika}). Asemblajo ek la prego-texti qui,}
\l{\cel{3}en la religio katolika, esas recitenda da la kleriki ye ula}
\l{\cel{3}kloki di la jorno. - DEFIS.}
\l{}
\l{\sou{brezo}. Ligno transformita a karbono per kombusto, sive ardor-}
\l{\cel{3}anta, sive extingita. - EFIS.}
\l{}
\l{\sou{brido}. La du rimeni qui, fixigita ye singla lateri di la}
\l{\cel{3}morso, uzesas por haltigar o direktar sel-kavalo o vetur-}
\l{\cel{3}kavalo. - EFIS.}
\l{}
\l{\sou{brigo}. (\sou{navig.})\cel{2}Navo qua havas nur un granda masto, inklinita}
\l{\cel{3}addope ad un mezeno-masto. - DEFIRS.}
\l{}
\l{\sou{brigado}. Korpo de armeani mi-diviziona. - DEFIRS.}
\l{}
\l{\sou{brigantino}. (\sou{navig.}) Mikra navo, rigoza quale brigo, ma sen-}
\l{\cel{3}ponta. - DEFIRS.}
}
